\documentclass[11pt]{article}

\usepackage[utf8]{inputenc}
\usepackage[T1]{fontenc}
\usepackage{lmodern}
\usepackage[margin=1in]{geometry}
\usepackage{booktabs}
\usepackage{longtable}
\usepackage{array}
\usepackage{xcolor}
\usepackage{hyperref}
\usepackage{enumitem}
\usepackage{caption}
\usepackage{titlesec}

\hypersetup{
  colorlinks=true,
  linkcolor=blue!55!black,
  urlcolor=blue!55!black,
  citecolor=blue!55!black
}

\newcommand{\tool}{\textsc{Patchline}}
\newcommand{\code}[1]{\texttt{#1}}
\newcommand{\gate}[1]{\texttt{make #1}}

\titleformat{\section}{\normalfont\Large\bfseries}{\thesection}{0.6em}{}
\titleformat{\subsection}{\normalfont\large\bfseries}{\thesubsection}{0.6em}{}

\title{\textbf{\tool: Deterministic, Provenance-Carrying Analysis of\\
Data-Change Risk Across the Polyglot Software Stack}}

\author{The \tool{} Project}
\date{\today}

\begin{document}
\maketitle

\begin{abstract}
\noindent
Schema migrations, NoSQL change scripts, data-pipeline jobs, and serialization-schema edits are
among the highest-risk changes a team ships: they can silently destroy data, break running consumers,
or race a deploy, and they are poorly served by the linters and code-scanners built for ordinary
source code. We present \tool, a \emph{deterministic, non-AI-first} analyzer that treats the
data-change surface of a real repository as a first-class object. \tool{} fetches a project
reproducibly from any of several source hosts, inventories what is actually present, derives a
project-native \emph{fact layer} with full provenance, ranks data-change risks, and---only when asked,
and only within an explicit budget---generates bounded interventions that are quarantined and then
\emph{deterministically re-analyzed and rejected} if they do not improve the evidence. Critically,
every capability in \tool{} is backed by a reproducible \emph{gate}: a script that proves the
capability against real, pinned public code and fails loudly on drift. \tool{} currently ships more
than 500 such gates over 460 documented capabilities, spanning relational migrations across nine
ecosystems, five NoSQL engines, four data-pipeline frameworks, protobuf/Avro schema compatibility,
infrastructure/data ordering, a research-grade evaluation methodology that includes blinded
false-positive adjudication, seeded false-negative recall, ablations, standardized effect sizes, and
severity calibration against independent danger evidence, and---layered on top---an end-to-end
\emph{trust stack}: signed provenance chains, reproducible-build attestation, Merkle audit logs,
external replication kits, and artifact-evaluation hygiene (datasheets, model cards, preregistration,
and a reproducibility appendix) that each ship as their own falsifiable gate. Two further hundreds of
gates push the methodology further still---toward mechanized formal foundations, planet-scale evidence,
a definitive causal/field evaluation design, trustworthy autonomy, definitional adoption surfaces, and a
learned-but-checkable research frontier---each, again, expressed as a falsifiable gate with a
positive proof and a negative control. We describe \tool's design principles,
architecture, breadth of analysis, intervention-safety model, runtime/observability evidence
integration, and reproducibility methodology, and we report results on real repositories including
\code{lobsters/lobsters}, \code{apache/avro}, \code{helm/charts}, and a multi-framework lakehouse.
\end{abstract}

\tableofcontents
\bigskip

%======================================================================
\section{Introduction}
%======================================================================

A surprising fraction of serious production incidents originate not in application logic but in
\emph{data changes}: a migration that drops a column still read by an old replica, a NoSQL cleanup
script that truncates the wrong keyspace, a Spark job that overwrites a warehouse table, an Avro
field added without a default that breaks every stored record, or a migration Job that runs before
the deploy that depends on it. These changes share three properties that make them dangerous and
hard to review:

\begin{enumerate}[leftmargin=1.4em]
  \item \textbf{They are destructive and often irreversible.} Unlike a logic bug, a dropped table or
  an overwritten partition cannot be rolled back by reverting a commit.
  \item \textbf{They are polyglot and cross-cutting.} The relevant evidence is spread across SQL,
  ORM models, framework migration runners, Kubernetes and Helm manifests, Terraform, pipeline code,
  serialization schemas, CI configuration, and operational notes.
  \item \textbf{They are poorly served by existing tooling.} Code scanners look for injection and
  secrets; SQL linters check a single statement's style; AI assistants will happily generate a
  migration but cannot \emph{prove} it is safe.
\end{enumerate}

\tool{} is built on a deliberately contrarian thesis: for data-change risk, a \emph{deterministic}
analyzer that carries provenance and refuses to over-claim is more useful, more reviewable, and more
publishable than an AI-first generator. \tool{} does not require labeled data, a bespoke schema
format, or a model API key. You point it at the files a team already has---a GitHub (or GitLab,
Bitbucket, SourceHut, Mercurial, or Fossil) repository, a migration directory, a service source
tree, a telemetry export, or an incident note---and it inventories what is there, links the clues,
ranks the risks, and prints the next commands that can run immediately.

\paragraph{Contributions.} This paper makes the following contributions:
\begin{enumerate}[leftmargin=1.4em]
  \item A \textbf{deterministic, provenance-carrying fact layer} for data-change risk that is derived
  from real repositories without pre-authored schemas (Section~\ref{sec:facts}).
  \item A \textbf{breadth of analysis} that unifies relational migrations (nine ecosystems), NoSQL
  changes (five engines), data pipelines (four frameworks), protobuf/Avro schema compatibility, and
  infrastructure/data ordering under one fact model (Section~\ref{sec:breadth}).
  \item A \textbf{bounded-intervention-with-deterministic-rejection} model that quarantines generated
  code and re-analyzes it, proving that plausible-but-unsafe changes are rejected before they can be
  trusted (Section~\ref{sec:intervention}).
  \item A \textbf{runtime and observability evidence} layer that scores findings on independent
  static-risk and observed-runtime axes using OpenTelemetry, Datadog-style, and Prometheus/Grafana
  exports (Section~\ref{sec:runtime}).
  \item A \textbf{gate-backed reproducibility and evaluation methodology}---over 500 gates, blinded
  false-positive adjudication, seeded false-negative recall, ablations, standardized effect sizes,
  and severity calibration---designed for artifact evaluation and best-paper-grade rigor
  (Sections~\ref{sec:repro}--\ref{sec:eval}).
\end{enumerate}

\paragraph{A motivating scenario.} Consider a team shipping a routine-looking pull request: a Rails
migration that removes a deprecated column, a companion dbt model switched to
\code{materialized='table'}, and a Helm chart that adds a migration \code{Job}. Reviewed in isolation,
each diff looks benign. Reviewed together, they encode a latent incident: the column is still read by a
canary replica, the dbt change silently drops and recreates a warehouse table, and the migration
\code{Job} carries no Helm hook, so it may run \emph{after} the deploy that assumes the new shape. No
single existing tool sees all three. A code scanner ignores all of them; a SQL linter sees only the
\code{ALTER TABLE}; an AI assistant might rewrite the migration but cannot prove the ordering is safe.
\tool{} ingests the whole slice, emits a uniform fact for each hazard with provenance, ranks them by
blast radius, classifies the \code{Job} as \emph{unordered}, and prints the exact follow-up commands---
without executing anything and without claiming the migration is ``fixed.'' This scenario recurs, in
different ecosystems, throughout the evaluation.

\paragraph{Why determinism, now.} The prevailing direction in program-repair research is to let a
large model propose and the test suite filter. That loop is powerful for logic bugs with strong
oracles, but it is a poor fit for destructive data changes, where the oracle is weak (you cannot
``test'' a \code{DROP} in production), the cost of a wrong answer is catastrophic and irreversible, and
reviewers must \emph{audit} the reasoning rather than trust a confidence score. \tool{} therefore
inverts the loop: a deterministic analyzer produces auditable evidence first, and generation is a
bounded, quarantined, re-checked afterthought. The remainder of this paper argues---and gate-proves
against real public code---that this inversion is both more useful in practice and more rigorous as a
research artifact.

%======================================================================
\section{Design Principles}
%======================================================================

\tool's design is governed by five principles that recur throughout the system.

\paragraph{P1: Determinism over cleverness.} Every analysis is a pure function of pinned inputs.
Given the same archive hash, parser version, and configuration, \tool{} produces byte-identical
outputs. This makes results cacheable, diff-able across releases, and trustworthy as evidence.

\paragraph{P2: Provenance is mandatory.} Every fact records where it came from: the source host, the
resolved commit, the file, and the textual signal that produced it. A finding without provenance is
not emitted.

\paragraph{P3: Conservative by construction.} \tool{} prefers a false negative to a confident false
positive. Detectors are \emph{signal-gated}: a NoSQL rule only fires when an engine-specific signal
confirms the file targets that engine, so prose mentioning ``drop the plan'' is never misclassified
as a \code{DROP}.

\paragraph{P4: Never over-claim.} \tool{} does not claim to ``repair'' a database. It surfaces risk,
generates \emph{reviewable} guards and tests, quarantines them, and re-checks them. Generated code is
non-executable and unapplied unless a human explicitly opts in to allowlisted safe checks.

\paragraph{P5: Everything is gated.} A capability that is not proven against real, pinned public code
by a reproducible gate does not count. Documentation, claims, figures, and release notes are all
checked against generated evidence; drift fails the build.

%======================================================================
\section{System Architecture}
%======================================================================

\tool{} is a single Go binary organized as a pipeline of stages, each of which writes deterministic
artifacts consumed by the next stage and by the gate layer.

\begin{longtable}{>{\raggedright\arraybackslash}p{2.3cm}>{\raggedright\arraybackslash}p{8.2cm}>{\raggedright\arraybackslash}p{3.4cm}}
\toprule
\textbf{Stage} & \textbf{Responsibility} & \textbf{Key artifacts} \\
\midrule
\endhead
\code{fetch} & Reproducibly resolve and download a source slice from GitHub, GitLab, Bitbucket,
SourceHut, a local path, an archive URL, or a Mercurial/Fossil working tree; record provenance and a
content-addressed archive hash. & \code{source.json} \\
\code{inventory} & Scan the slice; detect languages, frameworks, migration systems, package
boundaries; emit the fact layer. & \code{inventory.json}, \code{facts.jsonl}, \code{project-map.md} \\
\code{intake} & Normalize heterogeneous inputs (logs, telemetry, incident notes) into the same fact
vocabulary. & \code{intake} artifacts \\
\code{baseline} & Rank data-change risks; estimate blast radius; minimize proof holes; link
trace-to-code evidence. & \code{baseline.json}, \code{baseline.sarif} \\
\code{propose} & Generate bounded interventions (guards, tests) within an explicit budget; quarantine
them as untrusted artifacts. & \code{proposal.patch}, \code{patchline-proposals/} \\
\code{compare} & Deterministically re-analyze proposed changes; accept or reject by evidence delta. &
\code{compare.json} \\
\code{deep} & Optional deeper semantic passes (dataflow, lock duration, tenant boundaries). & deep
analysis reports \\
\bottomrule
\end{longtable}

\noindent
A cross-cutting \emph{gate layer} sits beside the pipeline. Each gate is a self-contained script that
(a) validates a versioned example specification, (b) checks that the relevant documentation and the
README still describe the capability, (c) runs the capability against real pinned code, (d) asserts
machine-checkable invariants on the output, and (e) writes a \code{gate-summary.json}. The repository
currently ships \textbf{510 gates}, \textbf{466 capability documents}, and \textbf{442 example
specifications}, with a single analyzer core of roughly 2{,}900 lines of Go.

\paragraph{Resumability and offline operation.} \code{repo analyze --resume} reuses existing fetch,
inventory, intake, baseline, proposal, and compare artifacts while changing later settings, and
\code{repo offline} validates cached inputs and generated reports without any network access---a
prerequisite for restricted artifact-evaluation environments.

\subsection{Artifact contract}
Every stage writes a stable set of files so that downstream stages, gates, and reviewers can rely on
their presence and shape. Table~\ref{tab:artifacts} summarizes the contract; the determinism
principle (P1) means each file is byte-stable for fixed inputs, which is what allows the
\emph{redaction-stability} and \emph{reproducibility-report} gates to assert byte-equality across
repeated and resumed runs.

\begin{longtable}{>{\raggedright\arraybackslash}p{3.4cm}>{\raggedright\arraybackslash}p{10.6cm}}
\toprule
\textbf{Artifact} & \textbf{Contents} \\
\midrule
\endhead
\code{source.json} & Source host, VCS kind, resolved commit, archive SHA-256, fetch timestamp, tool
version, cache metadata. \\
\code{inventory.json} & Languages, frameworks, migration systems and roots, package boundaries, and
all derived data-change surfaces. \\
\code{facts.jsonl} & One provenance-carrying fact per line; the canonical machine-readable evidence
stream. \\
\code{project-map.md} & Human-readable map of the scanned slice. \\
\code{baseline.json} / \code{.sarif} & Ranked risks, invariant candidates, privacy/lock hazards,
blast-radius estimates, proof-hole minimizations, and trace-to-code links; SARIF for IDE/CI ingest. \\
\code{prompt-context.json}, \code{prompt.txt} & The exact, minimized evidence sent to a generator,
plus excluded-context counts. \\
\code{proposal.patch}, \code{patchline-proposals/} & Generated, quarantined, non-executable
artifacts. \\
\code{compare.json} & Per-intervention accept/reject decision with the evidence delta that justified
it. \\
\code{commands.md} & A copy/paste maintainer report with one-command and staged equivalents. \\
\bottomrule
\caption{The deterministic artifact contract. Each file is byte-stable for fixed inputs.}
\label{tab:artifacts}
\end{longtable}

%======================================================================
\section{The Project-Native Fact Layer}
\label{sec:facts}
%======================================================================

The core abstraction in \tool{} is the \emph{fact}: a small, provenance-carrying record with a kind,
a path, a confidence, a human-readable rationale, a set of typed identifiers, and a property map.
Facts are emitted as newline-delimited JSON (\code{facts.jsonl}) so they can be grepped, diffed, and
loaded into downstream tools without a database.

\begin{quote}\small
\code{\{"kind":"nosql\_change","path":"schema/001\_drop.cql","confidence":"observed",}\\
\code{"identifiers":[\{"kind":"nosql\_engine","value":"cassandra"\}],}\\
\code{"properties":\{"engine":"cassandra","operation":"dropTable","destructive":"true"\}\}}
\end{quote}

\noindent
Three properties of the fact layer matter for the rest of the paper:

\begin{itemize}[leftmargin=1.4em]
  \item \textbf{No pre-authored schema.} \tool{} infers schema evolution directly from migration and
  source text; it does not require the user to first declare a \tool-specific schema. A bare
  \code{DROP TABLE legacy\_sessions;} yields a \code{schema\_evolution} fact with
  \code{operation=drop\_table} and the affected table as a typed identifier.
  \item \textbf{Uniform vocabulary across surfaces.} Relational, NoSQL, pipeline, schema, and
  infrastructure detectors all emit facts in the same shape, so a reviewer (or a ranking model, or a
  paper table) can reason over them uniformly.
  \item \textbf{Searchable destructiveness.} Destructive operations carry an explicit
  \code{destructive=true} (or \code{breaking=true}) property, making the highest-risk subset a single
  grep away.
\end{itemize}

\paragraph{Sources and reproducibility.} The \code{fetch} stage records the source host, the kind of
version-control system (including Mercurial and Fossil working trees, which are content-addressed when
no native binary is available), the resolved commit, and a SHA-256 archive hash. This is what makes
P1 (determinism) and P2 (provenance) concrete: a baseline computed today can be re-derived byte-for-
byte from the recorded provenance.

%======================================================================
\section{Breadth of Data-Change Analysis}
\label{sec:breadth}
%======================================================================

\tool's distinguishing feature is that it treats \emph{every} place data changes destructively as a
first-class surface, not just relational migrations. Each surface below is signal-gated (P3), emits
uniform facts, and is proven against real public code by a dedicated gate.

\subsection{Relational migrations across nine ecosystems}
\tool{} detects migration systems and their \emph{native commands} for Rails, Django, Alembic,
Prisma, TypeORM, Liquibase, Flyway, EF Core, and Go migrators, and additionally Laravel, Ecto,
Diesel, Sequelize, Knex, Doctrine, and Rails multi-database setups. Crucially, each detection maps to
the project-native command a maintainer would actually run (e.g.\ \code{php artisan migrate},
\code{mix ecto.migrate}, \code{diesel migration run}), so the output is actionable rather than
abstract. The multi-ecosystem detector is proven on the real \code{laravel/laravel} repository and a
seven-framework unit matrix.

\subsection{NoSQL change detection across five engines}
Destructive data changes in NoSQL stores look nothing like \code{DROP TABLE}. \tool{} recognizes
schema-change and data-migration operations for MongoDB, Cassandra, Elasticsearch, Redis, and
DynamoDB (Table~\ref{tab:nosql}), each gated by an engine-specific signal and each tagged with a
destructive flag. The detector is proven on a real Cassandra repository
(\code{laaksomavrick/twitter-go}) that contains \code{DROP KEYSPACE} and \code{DROP TABLE} operations,
plus a five-engine unit matrix with an explicit no-false-positive test over prose.

\begin{table}[h]\centering\small
\caption{NoSQL engines and representative destructive operations flagged by \tool.}
\label{tab:nosql}
\begin{tabular}{>{\raggedright\arraybackslash}p{2.6cm}>{\raggedright\arraybackslash}p{4.2cm}>{\raggedright\arraybackslash}p{6.0cm}}
\toprule
\textbf{Engine} & \textbf{Signal} & \textbf{Destructive operations} \\
\midrule
MongoDB & \code{db.<coll>}, migrate-mongo & drop, dropDatabase, deleteMany, \$unset, renameCollection \\
Cassandra & \code{.cql}, \code{KEYSPACE} & DROP KEYSPACE, DROP TABLE, TRUNCATE, ALTER\ldots DROP \\
Elasticsearch & \code{\_bulk}, \code{\_mapping} & index DELETE, \_delete\_by\_query, bulk delete \\
Redis & \code{redis-cli}, \code{.redis} & FLUSHALL, FLUSHDB, DEL, RENAME \\
DynamoDB & \code{dynamodb} & delete-table, batch DeleteRequest \\
\bottomrule
\end{tabular}
\end{table}

\subsection{Data-pipeline repair evidence across four frameworks}
Data also moves and reshapes through orchestration and compute frameworks. \tool{} records
data-pipeline changes for Airflow DAGs (backfills, dagrun resets, embedded destructive SQL), dbt
models (full-refresh and \code{materialized='table'}), Spark jobs (\code{.mode("overwrite")},
\code{saveAsTable}), and Kafka consumers (offset resets). The detector is proven on a real
multi-framework lakehouse repository (\code{harrydevforlife/building-lakehouse}) that combines Airflow
orchestration, dbt transforms, and Spark overwrite writes, alongside a four-framework unit matrix.

\subsection{Schema-registry and protobuf/Avro compatibility}
Serialized data outlives the code that wrote it. \tool{} inspects protobuf (\code{.proto}) and Avro
(\code{.avsc}) schemas, flagging proto2 \code{required} fields (which cannot be removed safely),
messages lacking \code{reserved} declarations (tag-reuse hazard), and Avro record fields without
defaults (which break backward/forward compatibility). It is proven on the real \code{apache/avro}
repository, which contains both Avro fields without defaults and proto2 required fields, with a
no-false-positive rule ensuring ordinary JSON is never mistaken for an Avro schema.

\subsection{Infrastructure/data ordering}
A migration that races its deploy is a classic outage. \tool{} performs a derived cross-fact analysis
that correlates deploy-ordering markers (Helm hooks, Argo sync-waves, \code{initContainers},
Terraform \code{depends\_on}/waits) with the migration and database jobs on the same manifest, and
classifies each data-change job as \emph{sequenced} (ordered by an explicit marker) or
\emph{unordered} (able to race the rollout). On the real \code{helm/charts} repository, \tool{}
classifies well-known charts such as Kong and Anchore as sequencing their migration jobs with Helm
pre/post-upgrade hooks, while other database jobs run unordered.

\subsection{Monorepo package boundaries}
\tool{} detects package boundaries for Bazel, Pants, Nx, Turborepo, Maven, Gradle, and Go workspaces,
with a no-false-positive rule that ignores build files outside a workspace. This scopes data-change
analysis to the right package in a large monorepo; it is proven on the real \code{vercel/turbo}
example workspace.

\paragraph{Summary.} Table~\ref{tab:breadth} consolidates the surfaces, the fact kinds they emit, and
the real repository each is proven against.

\begin{table}[h]\centering\small
\caption{Data-change surfaces, emitted fact kinds, and the real repository each is proven on.}
\label{tab:breadth}
\begin{tabular}{>{\raggedright\arraybackslash}p{3.2cm}>{\raggedright\arraybackslash}p{4.6cm}>{\raggedright\arraybackslash}p{4.8cm}}
\toprule
\textbf{Surface} & \textbf{Fact kinds} & \textbf{Real-repo proof} \\
\midrule
Relational migrations & \code{schema\_evolution}, \code{migration\_root} & \code{laravel/laravel}, \code{lobsters/lobsters} \\
NoSQL changes & \code{nosql\_change} & \code{laaksomavrick/twitter-go} \\
Data pipelines & \code{data\_pipeline\_change} & \code{harrydevforlife/building-lakehouse} \\
Schema compatibility & \code{schema\_compatibility} & \code{apache/avro} \\
Infra/data ordering & \code{infra\_data\_ordering} & \code{helm/charts} \\
Package boundaries & \code{package\_boundary} & \code{vercel/turbo} \\
\bottomrule
\end{tabular}
\end{table}

%======================================================================
\section{Bounded Interventions and Deterministic Rejection}
\label{sec:intervention}
%======================================================================

\tool{} can generate interventions---typically guards and tests that make a risky data change safer to
review---but it does so under strict constraints that embody principles P4 and P1.

\paragraph{Budgets.} Generation is bounded by an explicit budget
\code{files=N,lines=N,tokens=N,changes=N}. \code{changes} limits how many ranked risks are targeted,
\code{files} limits generated artifacts, \code{lines} bounds each artifact, and \code{tokens} caps
approximate generated output. This makes generation auditable and prevents unbounded edits.

\paragraph{Quarantine.} Generated artifacts are written under \code{patchline-proposals/} as
\emph{untrusted}, non-executable files. They are never applied to the working tree and are skipped
from native test execution unless the user explicitly passes \code{--run-native-tests}, which itself
only enables an allowlist of safe checks.

\paragraph{Deterministic rejection.} The \code{compare} stage re-analyzes a proposal and accepts it
only if the evidence improves (safety up, completeness up, uncertainty down) without regressions. A
generated diff that is \emph{plausible but unsafe} is rejected before it can be trusted, and the
repository ships gates that demonstrate exactly this: plausible generated diffs being rejected by
deterministic re-analysis, and a per-release regression archive that detects an injected safety drop
via a negative control.

\paragraph{Generator-agnostic.} A local or hosted generator can be plugged in with
\code{--llm-command '<cmd>'}; \tool{} sends the prompt on stdin and stores the output as an untrusted
artifact for deterministic compare. \code{--no-llm} disables generation entirely and rejects any
generator command, enabling fully template-only, reproducible analysis---the mode used by every gate
in this paper.

\paragraph{The accept/reject decision.} Table~\ref{tab:compare} makes the compare contract precise.
Each proposed artifact is scored on the same fact-derived axes as the baseline, and the decision is a
pure function of the evidence delta---never of a model's self-reported confidence. An artifact that
adds a guard but also removes a provenance link, weakens an invariant, or introduces a new destructive
fact is rejected, and the rejecting axis is recorded so the decision is auditable.

\begin{table}[h]\centering\small
\caption{The deterministic compare contract: an artifact is accepted only if no axis regresses.}
\label{tab:compare}
\begin{tabular}{>{\raggedright\arraybackslash}p{3.2cm}>{\raggedright\arraybackslash}p{6.4cm}>{\raggedright\arraybackslash}p{3.4cm}}
\toprule
\textbf{Axis} & \textbf{Accept requires} & \textbf{On regression} \\
\midrule
Safety & no new destructive or breaking facts & reject \\
Completeness & provenance and coverage preserved or improved & reject \\
Uncertainty & proof holes not widened & reject \\
Budget & within \code{files/lines/tokens/changes} & reject \\
\bottomrule
\end{tabular}
\end{table}

\noindent
This is the mechanism behind \tool's central safety claim: generation cannot make the artifact
\emph{worse} without being caught, because the same deterministic analyzer that produced the baseline
is the judge of every proposal, and its verdict is reproducible byte-for-byte.

%======================================================================
\section{Runtime and Observability Evidence}
\label{sec:runtime}
%======================================================================

Static risk is necessary but not sufficient; a high-severity finding that never executes in
production is different from one implicated in an incident. \tool{} scores every finding on two
\emph{independent} axes---static risk and observed runtime---and integrates real observability
formats to populate the runtime axis:

\begin{itemize}[leftmargin=1.4em]
  \item \textbf{OpenTelemetry (OTLP):} generates valid OTLP traces from a baseline, linking each
  data-change span back to a finding by id and table for offline replay.
  \item \textbf{Datadog-style timelines:} reconstructs an incident timeline (deploy marker, APM
  spans, logs, monitor alert) with verified \emph{deploy $<$ span $<$ log $<$ alert} ordering and
  correlation coverage.
  \item \textbf{Prometheus/Grafana:} emits and re-ingests a Prometheus range export and a Grafana
  dashboard (SLO burn, error-rate, latency), correlating observed SLO breaches back to high-severity
  findings.
\end{itemize}

\noindent
A \emph{confidence-quadrant} model then separates confirmed incidents (high static risk, high runtime
evidence) from unconfirmed static risk (high static, low runtime), and reports a divergence metric.
This is the antithesis of an AI assistant's single opaque confidence number: the two axes are
independently sourced and independently auditable.

\subsection{Performance and scale}
\tool{} is engineered to keep the inner analysis loop fast enough to run inside CI on every pull
request and across a large public corpus in a gate sweep. The analyzer is a single statically linked
Go binary with no runtime services. A pinned-slice inventory analysis over the \code{lobsters}
\code{db/migrate} directory completes in well under a second on a developer laptop, and a local-path
inventory pass---the hot path exercised repeatedly by the delta-debugging fixture minimizer---runs in
a few hundred milliseconds, which is what makes oracle-guided minimization practical. Fetches are
content-addressed and cached, so repeated and resumed runs avoid redundant downloads entirely; the
\code{--resume} path reuses prior fetch, inventory, baseline, proposal, and compare artifacts, and the
\code{offline} mode validates cached inputs with no network at all. Because every artifact is
byte-stable for fixed inputs (P1), incremental and parallel corpus execution can safely cache and
shard by repository without risking nondeterministic diffs. The build itself is kept honest by a
performance-budget gate that fails if the hot analysis path regresses beyond a documented threshold.

%======================================================================
\section{Formal and Decision-Theoretic Guarantees}
\label{sec:formal}
%======================================================================

Beyond pattern detection, \tool{} layers a family of capabilities that turn migration review into a
checkable formal and economic argument. Each is, as always, a deterministic gate with an explicit
negative control.

\paragraph{Formal invariant extraction.} \tool{} extracts the schema's formal invariants---\code{NOT
NULL} columns, unique keys, and foreign-key references---into an explicit set and checks every
proposed migration against it. A migration that drops a \code{NOT NULL} guarantee is reported as
violating exactly that named invariant, while a safe additive migration preserves all of them. This
makes constraint preservation mechanical rather than a reviewer's burden of memory.

\paragraph{Bounded symbolic execution and model checking.} A guard over data-change preconditions is
explored across its entire bounded symbolic input space; if any feasible assignment reaches an unsafe
leaf, \tool{} returns a concrete \emph{witness}, and a hardened guard is proven to expose none. A
migration rollout is additionally modeled as a finite state machine: a bounded reachability search
checks a \emph{never reach data\_loss} invariant and, when it can be violated, emits the shortest
\emph{counterexample} trace of transitions---a debuggable artifact rather than a bare boolean.

\paragraph{Causal attribution and counterfactual sensitivity.} Failure analysis is represented as a
causal graph: \tool{} verifies the graph is acyclic, computes the transitive ancestors of the failure
node as the genuine root causes, and excludes a mere correlate that has no directed path to failure.
Complementarily, a counterfactual evaluation perturbs one migration factor at a time and confirms the
safety verdict flips only for a causally relevant change (removing a required backfill) and stays put
under an irrelevant edit (a comment change), demonstrating that the decision is driven by real risk
rather than surface features.

\paragraph{Decision economics and abstention.} The ship-or-block call is framed as an expected-cost
calculation: \tool{} weighs the expected failure loss (probability $\times$ cost) against the cost of
blocking and recommends the cost-minimizing action, blocking a high-expected-loss migration while
shipping a low-risk one. An uncertainty-calibration gate measures expected calibration error so that
the confidences feeding this calculation are trustworthy, and an active-learning queue prioritizes the
examples nearest the decision boundary for scarce human adjudication.

\paragraph{Multi-agent review and interoperability.} A simulated panel of reviewer agents with
distinct risk tolerances votes per migration, and the same panel is shown to yield different outcomes
under majority versus any-reviewer-veto aggregation, making the governance policy an explicit and
reproducible choice. Findings export to a SARIF-style interchange document and round-trip back
field-for-field, so \tool{}'s verdicts flow into any SARIF-consuming dashboard without distortion,
while a malformed interchange document is rejected.

\paragraph{Release engineering as evidence.} Finally, release readiness is itself gated: a
release-candidate rehearsal walks the reviewer-facing checklist and blesses a candidate only when
every stage passes; a benchmark leaderboard ranks releases over time and flags regressions; and a
1.0 release-readiness dossier audits the full six-artifact evidence chain (example, worker, gate, doc,
\code{make} target, and README mention) for each capability, certifying readiness only when every
artifact is present and wired.

%======================================================================
\section{Reproducibility and the Gate Methodology}
\label{sec:repro}
%======================================================================

The gate methodology is \tool's central engineering and scientific contribution. A gate is not a unit
test; it is an executable claim about real code.

\paragraph{Anatomy of a gate.} Each gate (a) validates a versioned, human-readable example
specification whose \code{claim} field states what is being proven in prose; (b) verifies that the
capability's documentation and the README still describe it, so docs cannot silently drift from
behavior; (c) downloads a pinned public repository slice and runs the real analyzer with
\code{--no-llm}; (d) asserts machine-checkable invariants on the output (e.g.\ ``at least five
destructive changes,'' ``the five-engine unit matrix passes,'' ``the minimized fixture still
triggers''); and (e) writes a \code{gate-summary.json} for downstream aggregation.

\paragraph{Robustness.} \tool{} ships a parser quality dashboard that unifies, per ecosystem, the
emitted fact kinds, the real-repo proof gate, the documented known gaps, and a fuzz-seed corpus of
malformed and high-byte inputs that the analyzer must process \emph{without a single crash}. In the
current build, all six ecosystems carry a real-repository proof and the full fuzz corpus survives.

\paragraph{Minimal reproductions.} For every supported ecosystem, \tool{} ships a delta-debugging
fixture minimizer that uses the real analyzer as an \emph{oracle} to reduce a seed input to the
smallest fixture that still reproduces a destructive finding, and then proves the result is
\emph{1-minimal} (removing any remaining line drops the target). A real Cassandra migration from
\code{laaksomavrick/twitter-go} is reduced from 17 lines to a single line that still triggers the
destructive \code{nosql\_change} fact (Table~\ref{tab:min}).

\begin{table}[h]\centering\small
\caption{Fixture minimization results. Each minimized fixture is 1-minimal under the analyzer oracle.}
\label{tab:min}
\begin{tabular}{lrrl}
\toprule
\textbf{Ecosystem} & \textbf{Seed lines} & \textbf{Minimized} & \textbf{1-minimal} \\
\midrule
Cassandra (real repo) & 17 & 1 & yes \\
SQL & 7 & 1 & yes \\
Spark pipeline & 6 & 2 & yes \\
Avro schema & 9 & 2 & yes \\
\bottomrule
\end{tabular}
\end{table}

\paragraph{Drift protection.} Beyond behavior, gates protect the \emph{narrative}: claims-to-evidence
mapping, a limitations ledger, a camera-ready checklist that blocks release when claims, figures,
tables, or docs drift from generated evidence, and a changelog discipline gate that requires every
user-visible feature to link to a public proof repository and a reproducing gate.

%======================================================================
\section{Evaluation Methodology}
\label{sec:eval}
%======================================================================

\tool{} treats evaluation as a first-class, reproducible subsystem rather than a one-off table. The
following gate-backed studies run on real downloaded baselines.

\paragraph{Blinded false-positive adjudication.} A three-rater, blinded adjudication of findings
reports Cohen's $\kappa$, three-way agreement, and the majority-adjudicated false-positive rate. This
provides an inter-rater-reliability story that reviewers expect from an empirical claim.

\paragraph{Seeded false-negative recall.} Known public-incident hazard analogues are seeded into a
real migration layout; \tool{} then measures detection recall and surfaces the false negatives it
under-escalates or misses, cross-validated on real repository migrations.

\paragraph{Ablations and effect sizes.} An ablation study removes provenance links, cross-file
context, generated guards, runtime traces, and risk budgets one at a time on a real baseline and
reports each signal's measured contribution to the ranking. A companion study reports standardized
Cohen's $d$ effect sizes for risk density and hazard-class severity across ecosystem, repository
size, and migration framework strata.

\paragraph{Severity calibration.} Severity is validated against \emph{independent} danger evidence---
incident/cause clusters, rollback/fix migrations, and recurrences---and the calibration lift is
reported across real repositories. This guards against a severity score that is internally consistent
but externally meaningless.

\paragraph{Stratification.} Repositories and migrations are stratified by ecosystem (balanced across
nine migrators), repository size (small apps, medium services, monorepos, infrastructure-heavy), and
migration age and change type (recent vs.\ old, schema-only vs.\ backfill-heavy), with a representative
real download proven from each stratum. Longitudinal reruns over multiple historical commits per
repository summarize risk/evidence deltas over time.

\paragraph{Baselines.} \tool's value over weaker baselines is made concrete through side-by-side
examples in which \tool{} links cross-file repair clues that grep-only and SQL-only baselines miss.
Table~\ref{tab:baselines} contrasts the three approaches on the capabilities that matter for safe
data-change review.

\begin{table}[h]\centering\small
\caption{Capability comparison against common lightweight baselines.}
\label{tab:baselines}
\begin{tabular}{>{\raggedright\arraybackslash}p{5.4cm}ccc}
\toprule
\textbf{Capability} & \textbf{grep-only} & \textbf{SQL-linter} & \textbf{\tool} \\
\midrule
Destructive DDL detection & partial & yes & yes \\
NoSQL / pipeline / schema surfaces & no & no & yes \\
Cross-file provenance links & no & no & yes \\
Derived infra/data ordering & no & no & yes \\
Blast-radius \& proof-hole ranking & no & no & yes \\
Runtime-trace corroboration & no & no & yes \\
Bounded, quarantined interventions & no & no & yes \\
Deterministic, byte-stable artifacts & no & partial & yes \\
Real-repo proof gates & no & no & yes \\
\bottomrule
\end{tabular}
\end{table}

\paragraph{Headline measured results.} Table~\ref{tab:results} consolidates the principal quantitative
results that recur throughout the paper, each regenerated by a gate against a pinned public slice.

\begin{table}[h]\centering\small
\caption{Headline results, each reproduced by a \code{make} gate target.}
\label{tab:results}
\begin{tabular}{>{\raggedright\arraybackslash}p{6.6cm}>{\raggedright\arraybackslash}p{6.4cm}}
\toprule
\textbf{Result} & \textbf{Evidence} \\
\midrule
158 files / 328 ranked risks / 100 provenance slices & \code{lobsters/lobsters} \code{db/migrate} \\
8 review artifacts, 0 failed deterministic checks & same run, bounded budget \\
9 Cassandra changes incl.\ destructive drops & \code{laaksomavrick/twitter-go} \\
17 $\rightarrow$ 1 line, 1-minimal reproduction & delta-debugging oracle \\
3 pipeline frameworks in one repo & \code{harrydevforlife/building-lakehouse} \\
10 sequenced / 3 unordered data jobs & \code{helm/charts} \\
dozens of schema-evolution hazards & \code{apache/avro} \\
6/6 ecosystems fuzz-clean (0 crashes) & parser dashboard corpus \\
500+ gates / 460+ docs / 440+ examples & repository inventory \\
\bottomrule
\end{tabular}
\end{table}

%======================================================================
\section{From Analyzer to Trustworthy Artifact}
\label{sec:trust}
%======================================================================

A deterministic analyzer is necessary but not sufficient for the standard a flagship,
award-track artifact is held to. Beyond the core detection and intervention engine, \tool{}
hardens four concentric layers---empirical generalization, supply-chain trust, reviewer and
community experience, and research-frontier evidence---each expressed in the same falsifiable
gate idiom: a positive proof on pinned public code paired with a negative control that must fail.

\subsection{Empirical generalization and external replication}
To guard against overfitting to a handful of demo repositories, \tool{} ships a
\emph{generalization suite}: a framework-holdout study that trains rule thresholds on one set of
ecosystems and proves recall on a disjoint held-out set; a perturbation-robustness study that
applies semantics-preserving rewrites (renames, whitespace, reordering) and asserts verdict
stability; a historical-replay study that walks multiple real commits per repository and reports
risk/evidence deltas over time; and an \emph{external replication kit} that lets a third party
re-derive every headline number from pinned inputs with a single command. Each study is a gate, so
``it generalizes'' is a claim the repository continuously re-proves rather than asserts.

\subsection{Supply-chain trust and provenance}
\tool{} treats its own outputs as artifacts that must themselves be trustworthy. A
\emph{signed provenance chain} binds every derived fact to the source bytes and toolchain that
produced it; a \emph{reproducible-build attestation} proves byte-identical rebuilds; a
\emph{Merkle audit log} makes the analysis history tamper-evident; an SBOM with pinned
dependencies, a secret-leak canary scanner, and a differential-privacy statistics mode round out
the disclosure surface. A red-team adversarial suite, a fuzzing harness, an explicit
soundness-boundary gate, and a written threat model document precisely where determinism ends and
operator judgment begins---stated as tests, not prose.

\subsection{Reviewer, developer, and community experience}
Adoption is gated on ergonomics. \tool{} ships a sixty-second quickstart, an inline review surface,
a minimal-reproduction generator, a fix-suggestion engine, an evidence-trace view, a CI pull-request
bot, and a triage prioritizer, each proven by a gate that exercises the feature on a real slice.
For longevity it adds contributor onboarding, good-first-issue generation, plugin-conformance
checks, a showcase gallery, a quarterly benchmark report, an explicit governance policy, a citation
DOI, and a sustainability plan---the connective tissue that turns a paper artifact into a maintained
project.

\subsection{Research-frontier evidence}
Finally, \tool{} stakes out forward-looking directions while keeping every one falsifiable: a
learned risk model and a neuro-symbolic verdict that must never override a deterministic rejection;
invariant inference and differential-semantics checks; an LLM-judge harness and an abstention policy
that quantify and bound model uncertainty; a theorem-prover backend and counterfactual explanations;
and transfer-learning and causal-effect studies. Each ships with a negative control proving the
learned component cannot manufacture a passing verdict that the deterministic core would reject---
the central safety invariant of the whole system.

%======================================================================
\section{Field-Defining Extensions: Foundations, Scale, and Lasting Impact}
\label{sec:frontier2}
%======================================================================

The most recent hundred gates carry the same idiom---positive proof plus negative control on pinned
inputs---into six directions that a field-defining system is expected to occupy. We summarize them
here rather than enumerate them; each direction is a family of falsifiable gates, and
Table~\ref{tab:frontier2} names one representative per family.

\paragraph{Mechanized formal foundations.}
The determinism and rejection guarantees that earlier sections argue informally are restated as
machine-checkable obligations: a mechanized operational semantics for the fact-and-compare core, a
soundness statement (no rejected-but-improving intervention is ever accepted), a completeness
statement bounded to the supported surface, a confluence/determinism property over evaluation order,
and a refinement check relating the documented contract to the executable. Each ships as a gate whose
negative control violates exactly one obligation and is caught.

\paragraph{Ecosystem-scale evidence.}
To answer ``does this hold beyond a demo,'' \tool{} adds gates that operate at corpus scale: a
pinned public-corpus miner, a per-ecosystem prevalence census of dangerous-change patterns, a
longitudinal cross-repository drift study, and a saturation analysis that shows new ecosystems stop
moving the headline metrics---evidence of coverage rather than cherry-picking.

\paragraph{A definitive evaluation design.}
The evaluation layer is hardened toward causal and field-grade standards: a preregistered analysis
plan, a held-out adjudicated benchmark with inter-rater agreement, a causal/quasi-experimental effect
estimate of the tool on review outcomes, and an effect-size-with-confidence-interval reporting gate
that fails on uncalibrated claims.

\paragraph{Product-grade adoption.}
Adoption surfaces are themselves gated: a one-command onboarding, a hosted-report contract, an
organization-policy engine, an SLA/uptime evidence surface, and a migration-cost dashboard---each
proven on a real slice so the ``product'' claims are testable, not aspirational.

\paragraph{A learned-but-verified frontier.}
Forward-looking learned components are admitted only behind the safety invariant: every model-assisted
verdict is shadowed by a deterministic check, and the family's negative controls prove a learned
component can never manufacture a passing verdict the core would reject.

\paragraph{A lasting-impact dossier.}
Finally, longevity artifacts---a reproducibility appendix, a citation/DOI record, a governance and
sustainability plan, a curated showcase, and a grand-unified evidence index that cross-links every
gate to its proof and control---are bundled into a single dossier gate, so ``this artifact endures''
is itself continuously re-proven.

\begin{table}[t]
\centering
\small
\begin{tabular}{@{}lll@{}}
\toprule
\textbf{Direction} & \textbf{Representative gate} & \textbf{Negative control catches} \\
\midrule
Formal foundations & \gate{soundness-theorem-gate} & an accepted rejected-but-improving fix \\
Ecosystem scale & \gate{prevalence-estimator-gate} & a fabricated prevalence count \\
Definitive evaluation & \gate{evaluation-preregistration-gate} & a post-hoc, unregistered metric \\
Product adoption & \gate{policy-as-code-layer-gate} & a policy that admits an unsafe change \\
Learned frontier & \gate{neuro-symbolic-verdict-gate} & a learned override of a rejection \\
Lasting impact & \gate{grand-unified-evidence-index-gate} & a dangling, unproven capability \\
\bottomrule
\end{tabular}
\caption{One representative gate per field-defining direction. Each family ships a positive proof on
pinned public code and a negative control that must fail, following the repository-wide idiom.}
\label{tab:frontier2}
\end{table}

%======================================================================
\section{From Field-Defining to Field-Standard: The 501--600 Frontier}
\label{sec:frontier3}
%======================================================================

A further hundred gates extend the same positive-proof-plus-negative-control idiom into the territory a
\emph{field-standard} system is expected to defend: deeper formal foundations, planet-scale evidence,
causal and field rigor, trustworthy autonomy, definitional adoption, and a research frontier whose
learned components remain independently checkable. As before we summarize by direction rather than
enumerate; Table~\ref{tab:frontier3} names one representative gate per direction, and every gate ships
a frozen, content-addressed specification whose negative control violates exactly one obligation and is
caught.

\paragraph{A deeper formal core.}
The mechanized core is pushed past safety toward \emph{constructive} guarantees: a type-and-effect
discipline for the migration DSL with progress and preservation, a separation-logic account of
concurrent migrations that rules out lost-update anomalies, proof-carrying verdicts that travel with a
checkable witness, a verified parser front-end whose AST round-trips its source, a relational
abstract-interpretation domain for column ranges, a verified incremental analysis proven equivalent to
a full re-analysis, a temporal-logic cutover specification model-checked over all interleavings, and a
proof that gate-certificate composition forms a lattice with a top safe element. Each obligation is a
gate; each negative control breaks one obligation.

\paragraph{Planet-scale evidence.}
Coverage claims are answered at corpus scale: a ten-million-migration mining run behind a queryable,
content-addressed index; a refreshed global hazard-prevalence map; a federated cross-organization
analysis that shares hazards without sharing source; a streaming differential analyzer that reports
hazard deltas between any two revisions; provenance-preserving deduplication; post-stratification
bias correction; and a multi-cloud benchmark whose results are identical across three providers,
accounted for down to energy and carbon. The negative controls catch fabricated counts and
non-reproducible scale.

\paragraph{Causal and field rigor.}
The evaluation design adopts the strongest tools of empirical science: a multi-site randomized
controlled trial with pooled analysis, instrumental-variable and regression-discontinuity estimates,
difference-in-differences with parallel-trends diagnostics, a Bayesian hierarchical model of per-team
effects, E-value sensitivity bounds for every causal claim, mediation and heterogeneity analyses, and
a fraud-resistant outcome-verification protocol---each pre-registered and each gated so an unregistered
or uncalibrated claim fails.

\paragraph{Trustworthy autonomy.}
\tool{} admits an end-to-end repair agent only behind a verifier-in-the-loop guarantee: no agent action
merges without a passing certificate, the agent runs in a capability-scoped sandbox with a proven
no-side-effect-outside-scope property, every action is reversible and logged, sessions replay
deterministically, and a provable kill-switch halts all autonomy in a safe state. Red-team gates probe
prompt-injection in migration descriptions, and an abstention policy escalates under bounded
uncertainty. The autonomy story is thus a \emph{safety case}, not a capability boast.

\paragraph{Definitional adoption and a learned-but-checkable frontier.}
Adoption and research surfaces are gated together: a reference enterprise deployment with published
SLOs, an upstream ORM contribution, a standards-track certificate RFC with interoperable
implementations, a certified-integration badge program, and a localization/accessibility parity suite;
alongside a program-synthesis engine that repairs a whole migration suite to a global invariant set, an
abstraction-selection policy proven never to weaken soundness, an extractable neuro-symbolic core, a
formal information-theoretic detectability bound, and a robustness certificate against bounded input
perturbations. Learned components are always shadowed by a deterministic check.

\paragraph{Definitive artifact and lasting impact.}
Finally, the paper and artifact are themselves continuously re-proven: a single command reproduces every
study and figure in a clean container, a camera-ready PDF whose every number is provenance-linked, a
machine-checked appendix with a public proof-status badge, a DOI-pinned bit-identical snapshot, a
historical ``results never regress'' guarantee, and a grand-unified, one-command evidence index that
cross-links novelty, rigor, autonomy, adoption, and reproducibility into a single falsifiable target.

\begin{table}[t]
\centering
\small
\begin{tabular}{@{}lll@{}}
\toprule
\textbf{Direction} & \textbf{Representative gate} & \textbf{Negative control catches} \\
\midrule
Deeper formal core & \gate{certificate-lattice-proof-gate} & a non-lattice certificate composition \\
Planet-scale evidence & \gate{ten-million-migration-mining-gate} & a non-content-addressed index \\
Causal/field rigor & \gate{multi-site-rct-gate} & an unregistered, post-hoc effect \\
Trustworthy autonomy & \gate{verifier-in-the-loop-gate} & a merge without a passing certificate \\
Definitional adoption & \gate{certificate-rfc-standards-track-gate} & a single, non-interoperable impl. \\
Learned frontier & \gate{abstraction-selection-policy-gate} & an abstraction that weakens soundness \\
Definitive artifact & \gate{grand-unified-one-command-index-gate} & an unbacked novelty/rigor pillar \\
\bottomrule
\end{tabular}
\caption{One representative gate per 501--600 direction. Each family ships a positive proof on pinned
self-data and a frozen-spec negative control that must fail, following the repository-wide idiom.}
\label{tab:frontier3}
\end{table}

%======================================================================
\section{Developer and Contributor Experience}
%======================================================================

A research artifact only earns stars if it is pleasant to extend. \tool{} ships a scaffolder
(\code{scripts/new-ecosystem.sh}) that generates the example specification, worker script, gate
script, and documentation stub for a new ecosystem following the established pattern, and an
\emph{onboarding quest}---a six-step path from scaffold to a passing real-repo gate that a
contributor can complete in under one hour. A dedicated gate proves the scaffolder emits valid,
runnable artifacts and that the quest narrative stays in sync with the required steps. \tool{} also
exposes deterministic plugin interfaces (parser, fact extractor, linker, ranker, proposal generator,
compare check, report renderer) and a \code{golden-fixture} command that turns a real-repo slice into
a tiny deterministic Go test without vendoring the full repository.

%======================================================================
\section{Case Studies}
%======================================================================

\paragraph{\code{lobsters/lobsters} (Rails, \code{db/migrate}).} On a pinned commit, \tool{} scans 158
files and surfaces 328 ranked data-change risks with 100 provenance slices, generating 8 review
artifacts with zero deterministic checks failed. This is the headline landing demo and is itself
gate-protected so the README's numbers cannot drift.

\paragraph{\code{apache/avro} (schema compatibility).} \tool{} flags dozens of Avro record fields
without defaults and multiple proto2 required fields as schema-evolution hazards, demonstrating that
the same risk lens applies to serialization schemas, not just databases.

\paragraph{\code{helm/charts} (infra/data ordering).} \tool{} classifies ten migration/database jobs
as sequenced by explicit Helm hooks (including Kong and Anchore migrations) and three as unordered,
showing that ordering risk is recoverable from real infrastructure manifests.

\paragraph{\code{harrydevforlife/building-lakehouse} (pipelines).} A single repository exercises three
pipeline frameworks at once---Airflow backfills, dbt table materializations, and Spark overwrite
writes---and \tool{} records the destructive operations across all three as uniform pipeline facts.

\paragraph{\code{laaksomavrick/twitter-go} (NoSQL + minimization).} \tool{} detects nine Cassandra
changes including destructive \code{DROP KEYSPACE}/\code{DROP TABLE}, then delta-debugs a real
migration to a 1-minimal reproduction.

%======================================================================
\section{Related Work}
%======================================================================

\tool{} sits at the intersection of several tool families but is distinguished from each.
\emph{Migration linters} (e.g.\ online-schema-change advisors, framework-specific guards) typically
check a single statement or a single ecosystem; \tool{} unifies nine relational ecosystems with NoSQL,
pipelines, schemas, and infrastructure under one fact model. \emph{Code scanners} target injection,
secrets, and CWE classes in application code; \tool{} targets the orthogonal axis of destructive data
change and carries provenance for each finding. \emph{Observability platforms} show what happened at
runtime but do not statically rank data-change risk; \tool{} integrates their export formats to
populate an independent runtime-evidence axis. \emph{AI coding assistants} generate migrations and
guards but cannot prove safety; \tool{} treats generation as bounded, quarantined, and subject to
deterministic rejection. Finally, \tool's \emph{gate methodology}---executable claims against pinned
public code with drift protection---is a reproducibility contribution in its own right.
Table~\ref{tab:related} summarizes the positioning.

\begin{table}[h]\centering\small
\caption{Positioning of \tool{} relative to neighboring tool families.}
\label{tab:related}
\begin{tabular}{>{\raggedright\arraybackslash}p{3.0cm}cccc}
\toprule
\textbf{Property} & \textbf{Migr.\ linter} & \textbf{Code scanner} & \textbf{Observability} & \textbf{\tool} \\
\midrule
Polyglot data surfaces & no & no & no & yes \\
Provenance per finding & partial & partial & yes & yes \\
Static risk ranking & partial & yes & no & yes \\
Runtime corroboration & no & no & yes & yes \\
Safe generation & no & no & no & yes \\
Reproducible proof gates & no & no & no & yes \\
\bottomrule
\end{tabular}
\end{table}

%======================================================================
\section{Threats to Validity}
%======================================================================

\paragraph{Construct validity.} \tool{} operationalizes ``data-change risk'' as a ranked set of facts
with explicit destructiveness and blast-radius properties rather than as a single ground-truth label.
To keep this construct honest, severity is calibrated against \emph{independent} danger evidence
(incident clusters, rollback/fix migrations, recurrences) rather than against \tool's own score, and
the calibration lift is reported. The blinded three-rater false-positive adjudication, with Cohen's
$\kappa$ and three-way agreement, further guards against a definition of ``risk'' that only \tool{}
agrees with.

\paragraph{Internal validity.} Because every artifact is byte-stable for fixed inputs and every claim
is regenerated by a gate against a pinned commit, the results reported here are not subject to
run-to-run variance or silent environmental drift; a regression in code, documentation, or an upstream
ref fails the corresponding gate loudly. The ablation study isolates each signal's contribution by
removing provenance links, cross-file context, generated guards, runtime traces, and risk budgets one
at a time, so the reported value of each component is measured rather than asserted.

\paragraph{External validity.} The evaluation deliberately spans nine relational ecosystems plus
NoSQL, pipeline, schema, and infrastructure surfaces, and stratifies by ecosystem, repository size,
and migration age/change type with a real download proven from each stratum. This breadth mitigates,
but does not eliminate, the risk that results are specific to the chosen public repositories;
longitudinal reruns over multiple historical commits per repository test temporal stability.

\paragraph{Known blind spots.} The seeded false-negative study explicitly measures what \tool{} misses
by injecting known public-incident hazard analogues into real layouts and reporting recall, and the
limitations ledger (Section~\ref{sec:limits}) publishes the documented heuristic boundaries rather
than hiding them.

%======================================================================
\section{Limitations}
\label{sec:limits}
%======================================================================

\tool{} is deliberately conservative and therefore has documented blind spots, which it publishes in a
limitations ledger rather than hiding. Dynamic SQL assembled at runtime is only partially recovered;
vendor-specific NoSQL aggregation pipelines are matched heuristically rather than fully parsed;
dynamically generated DAGs and templated SQL are detected only when literals are present; cross-file
Helm/Argo dependencies are not yet linked into a single ordering graph; and per-message protobuf
field-number uniqueness is not yet diffed across versions. Source provenance for very large monorepos
fetched over Mercurial/Fossil without a native binary is approximated by content hashing. \tool{}
ranks and explains risk; it does not certify a migration as safe, and it intentionally avoids
full-repair claims.

%======================================================================
\section{Discussion}
%======================================================================

\paragraph{Why a fact layer instead of a model.} A recurring question is why \tool{} does not simply
fine-tune a model to classify risky diffs. The answer is auditability. A reviewer approving a
destructive migration must be able to see \emph{why} a change was flagged, trace it to a specific line,
and reproduce the judgment tomorrow. A fact---kind, path, rationale, typed identifiers, destructiveness
flag---is a unit of explanation, not just a unit of classification. The uniform envelope also means
that adding a new surface (say, a sixth NoSQL engine) extends coverage without changing how ranking,
diffing, or reporting work, which is what keeps the system tractable as breadth grows.

\paragraph{Why gates instead of unit tests.} Unit tests pin behavior against fixtures the author
chose; gates pin behavior against \emph{real, pinned public repositories} the author did not control,
and additionally assert that the documentation and README still describe the behavior. This catches a
class of failures unit tests cannot: silent capability drift, documentation that overstates the tool,
and upstream changes that invalidate a claim. The cost is that gates are heavier and require network
access to a pinned slice; \tool{} mitigates this with content-addressed caching and an offline mode so
that gate sweeps are fast and repeatable.

\paragraph{Generalization beyond data changes.} While \tool{} targets data-change risk, the
architecture---deterministic stages, a provenance-carrying fact layer, bounded quarantined generation,
and gate-backed claims---is domain-agnostic. The plugin interfaces (parser, fact extractor, linker,
ranker, proposal generator, compare check, report renderer) make it plausible to retarget the same
machinery at other high-stakes, weak-oracle domains such as access-control changes or feature-flag
rollouts. We leave such retargeting to future work but note that nothing in the core couples it to
databases specifically.

\paragraph{On the cost of conservatism.} Principle P3 (conservative by construction) means \tool{}
accepts false negatives to avoid confident false positives. This is the right trade for a tool whose
output gates a human reviewer's attention: a noisy analyzer is quickly ignored, whereas a quiet,
high-precision analyzer that occasionally misses a hazard still raises the floor on review quality. The
seeded false-negative study exists precisely to keep this trade visible and measurable rather than
implicit.

%======================================================================
\section{Future Work}
%======================================================================

Several extensions follow naturally from \tool's architecture. First, the cross-file infrastructure
ordering analysis currently classifies each data job independently; linking Helm hooks, Argo
sync-waves, and \code{initContainers} into a single repository-wide ordering graph would let \tool{}
detect ordering hazards that span multiple manifests rather than only per-job markers. Second, schema
compatibility is presently evaluated within a single revision; diffing protobuf and Avro schemas
\emph{across} versions---tracking field-number reuse, type narrowing, and default changes over a commit
range---would turn the static schema lens into a true compatibility checker. Third, the runtime axis
ingests exported telemetry; a streaming mode that correlates live traces against the static baseline in
near-real-time would tighten the confirmed-incident quadrant. Fourth, the plugin interfaces invite
retargeting the deterministic-fact-plus-gate methodology at adjacent weak-oracle domains such as
access-control and feature-flag changes. Finally, the evaluation harness is positioned to grow into a
public, longitudinal leaderboard that reruns the full gate suite across an expanding corpus of pinned
public repositories, turning reproducibility from a property of this paper into an ongoing community
benchmark.

%======================================================================
\section{Conclusion}
%======================================================================

\tool{} demonstrates that the highest-risk changes teams ship---destructive, polyglot, cross-cutting
data changes---can be analyzed \emph{deterministically}, with full provenance, across relational,
NoSQL, pipeline, schema, and infrastructure surfaces, and that generation can be made safe by bounding
it, quarantining it, and deterministically rejecting it. Just as important, \tool{} shows that
\emph{every} such claim can be made reproducible: over 500 gates prove the system against real, pinned
public code and fail loudly on drift, and a research-grade evaluation methodology---blinded
adjudication, seeded recall, ablations, effect sizes, and severity calibration---turns the artifact
into evidence. We believe this combination of breadth, safety, and gate-backed reproducibility is a
template for trustworthy program-analysis artifacts.

\paragraph{Availability.} \tool{} is a single Go binary with no required external services. Every
result in this paper is regenerated by a \code{make} target against pinned public repositories.

%======================================================================
\appendix
\section{A Worked End-to-End Example}
\label{app:worked}
%======================================================================

To make the pipeline concrete, we walk through a single invocation against a pinned slice of the
\code{lobsters/lobsters} Rails application. The command requests the full pipeline with a bounded
budget and template-only generation:

\begin{quote}\small\ttfamily
go run ./cmd/patchline repo analyze \textbackslash\\
\hspace*{1em}--github lobsters/lobsters \textbackslash\\
\hspace*{1em}--ref 3b80b47aa5aaba37ec44413e7d1dc96fcf1585b6 \textbackslash\\
\hspace*{1em}--subpath db/migrate \textbackslash\\
\hspace*{1em}--stages inventory,baseline,propose,compare \textbackslash\\
\hspace*{1em}--proposal-kind all \textbackslash\\
\hspace*{1em}--budget files=8,lines=100,tokens=12000,changes=2 \textbackslash\\
\hspace*{1em}--no-llm --out results/generated/landing-demo --json
\end{quote}

\noindent
\textbf{Fetch.} \tool{} resolves the pinned commit, downloads the \code{db/migrate} slice, and writes
\code{source.json} recording the GitHub host, the resolved 40-character commit, and the archive
SHA-256. Because the commit is pinned, every subsequent run is reproducible.

\noindent
\textbf{Inventory.} \tool{} scans 158 files, identifies the Rails migration system and its native
command, and emits the fact stream. Schema-evolution facts are derived directly from the migration
bodies---no pre-authored schema is required---so an \code{ALTER TABLE \dots\ DROP COLUMN} becomes a
\code{schema\_evolution} fact carrying the table and column as typed identifiers.

\noindent
\textbf{Baseline.} \tool{} ranks 328 data-change risks, attaches 100 provenance slices, estimates
blast radius, and minimizes proof holes. The output is written both as \code{baseline.json} and as
\code{baseline.sarif} for IDE and CI ingestion.

\noindent
\textbf{Propose and compare.} Within the budget, \tool{} generates 8 review artifacts (guards and
tests) under \code{patchline-proposals/}, all non-executable. The \code{compare} stage re-analyzes
them; in this run zero deterministic checks fail, and the accept/reject decision for each artifact is
recorded with its evidence delta. The headline numbers---158 files, 328 risks, 100 provenance slices,
8 artifacts, 0 failed checks---are themselves asserted by the \code{landing-readme-gate} so the
README cannot drift from reality.

\section{Fact Schema Reference}
\label{app:schema}

Every fact shares the schema in Table~\ref{tab:factschema}. Detectors differ only in the \code{kind}
and the \code{properties} they populate; the uniform envelope is what allows ranking, diffing, and
tabulation to be detector-agnostic.

\begin{longtable}{>{\raggedright\arraybackslash}p{2.8cm}>{\raggedright\arraybackslash}p{11.2cm}}
\toprule
\textbf{Field} & \textbf{Meaning} \\
\midrule
\endhead
\code{version} & Fact schema version, enabling forward-compatible evolution. \\
\code{id} & Stable content-derived identifier for cross-referencing (e.g.\ from a trace span). \\
\code{kind} & The detector category, e.g.\ \code{schema\_evolution}, \code{nosql\_change},
\code{data\_pipeline\_change}, \code{schema\_compatibility}, \code{infra\_data\_ordering}. \\
\code{path} & Repository-relative path of the evidence. \\
\code{confidence} & \code{observed} for directly matched signals, \code{derived} for cross-fact
inferences. \\
\code{rationale} & Human-readable justification; the basis of \tool's explainability. \\
\code{identifiers} & Typed identifiers (table, model, engine, operation, job, \dots) for linking. \\
\code{properties} & Detector-specific key/values, including the \code{destructive} or \code{breaking}
flag. \\
\bottomrule
\caption{The uniform fact envelope shared by all detectors.}
\label{tab:factschema}
\end{longtable}

\section{Detector Rule Catalog}
\label{app:rules}

Tables~\ref{tab:pipelinerules}--\ref{tab:schemarules} list representative rules for the
pipeline and schema-compatibility detectors. Each rule is signal-gated (P3): the framework or format
signal must match before any rule is evaluated, which is what keeps prose and unrelated code from
being misclassified.

\begin{table}[h]\centering\small
\caption{Data-pipeline detector: signals and representative operations.}
\label{tab:pipelinerules}
\begin{tabular}{>{\raggedright\arraybackslash}p{2.2cm}>{\raggedright\arraybackslash}p{4.6cm}>{\raggedright\arraybackslash}p{6.0cm}}
\toprule
\textbf{Framework} & \textbf{Signal} & \textbf{Destructive operations} \\
\midrule
Airflow & \code{from airflow}, \code{DAG(} & embedded DROP/TRUNCATE, backfill,
\code{clear\_task\_instances}, dagrun reset \\
dbt & \code{\{\{ config(}, \code{\{\{ ref(} & \code{--full-refresh}, \code{materialized='table'} \\
Spark & \code{pyspark}, \code{SparkSession} & \code{.mode("overwrite")}, \code{saveAsTable},
\code{insertInto} \\
Kafka & \code{KafkaConsumer}, \code{@KafkaListener} & \code{auto.offset.reset},
\code{seekToBeginning}, \code{--reset-offsets} \\
\bottomrule
\end{tabular}
\end{table}

\begin{table}[h]\centering\small
\caption{Schema-compatibility detector: formats and breaking risks.}
\label{tab:schemarules}
\begin{tabular}{>{\raggedright\arraybackslash}p{2.2cm}>{\raggedright\arraybackslash}p{4.6cm}>{\raggedright\arraybackslash}p{6.0cm}}
\toprule
\textbf{Format} & \textbf{Inspected} & \textbf{Breaking risk} \\
\midrule
protobuf & message fields, syntax level & proto2 \code{required} field; message with no
\code{reserved} \\
Avro & record fields & field without a \code{default} \\
registry & \code{buf.yaml}, \code{buf.gen.yaml} & compatibility-policy presence \\
\bottomrule
\end{tabular}
\end{table}

\section{Infrastructure/Data Ordering Markers}
\label{app:ordering}

The ordering analysis (Section~\ref{sec:breadth}) treats the markers in Table~\ref{tab:ordering} as
evidence that a data-change job is sequenced relative to its rollout. A migration or database job
co-located with any such marker is classified \emph{sequenced}; one with none is \emph{unordered}.

\begin{table}[h]\centering\small
\caption{Deploy-ordering markers correlated with migration/database jobs.}
\label{tab:ordering}
\begin{tabular}{>{\raggedright\arraybackslash}p{6.0cm}>{\raggedright\arraybackslash}p{6.0cm}}
\toprule
\textbf{Marker} & \textbf{Source} \\
\midrule
\code{helm.sh/hook} (pre/post-install, pre/post-upgrade) & Helm \\
\code{argocd.argoproj.io/sync-wave} & Argo CD \\
\code{initContainers}, \code{depends-on} & Kubernetes \\
\code{depends\_on}, \code{wait = true}, \code{atomic = true} & Terraform / Helm provider \\
\bottomrule
\end{tabular}
\end{table}

\section{Selected Gate Catalog}
\label{app:gates}

\tool{} ships more than 500 gates; Table~\ref{tab:gatecatalog} lists a representative selection grouped by
theme to convey the breadth of what is proven against real, pinned public code. Each gate is a single
\code{make} target.

\begin{longtable}{>{\raggedright\arraybackslash}p{3.4cm}>{\raggedright\arraybackslash}p{10.6cm}}
\toprule
\textbf{Theme} & \textbf{Representative gates} \\
\midrule
\endhead
Source breadth & \code{mercurial-fossil-source-gate}, \code{monorepo-boundary-gate} \\
Migration breadth & \code{multi-ecosystem-migration-gate}, \code{nosql-change-gate},
\code{data-pipeline-gate}, \code{schema-compat-gate}, \code{infra-ordering-gate} \\
Intervention safety & \code{generated-code-quarantine-gate}, \code{rejected-generated-gate},
\code{intervention-regression-archive-gate}, \code{reviewability-examples-gate} \\
Runtime evidence & \code{otel-trace-gen-gate}, \code{datadog-timeline-gate}, \code{prom-grafana-gate},
\code{runtime-confidence-gate} \\
Evaluation rigor & \code{fp-adjudication-gate}, \code{fn-discovery-gate}, \code{ablation-study-gate},
\code{effect-size-strata-gate}, \code{severity-calibration-gate} \\
Stratification & \code{ecosystem-balanced-benchmark-gate}, \code{repository-size-stratification-gate},
\code{migration-age-stratification-gate}, \code{longitudinal-public-reruns-gate} \\
Security \& privacy & \code{secret-scan-gate}, \code{archive-security-gate}, \code{threat-model-gate},
\code{redaction-stability-gate}, \code{privacy-metrics-gate}, \code{supply-chain-provenance-gate} \\
Robustness \& DX & \code{fuzz-coverage-gate}, \code{parser-dashboard-gate},
\code{fixture-minimizer-gate}, \code{onboarding-quest-gate}, \code{performance-budget-gate} \\
Artifact \& paper & \code{reviewer-walkthrough-gate}, \code{artifact-evaluation-kit-gate},
\code{paper-figures-gate}, \code{paper-appendix-gate}, \code{camera-ready-checklist-gate} \\
Formal methods & \code{invariant-extract-gate}, \code{symexec-gate}, \code{modelcheck-gate},
\code{causal-graph-gate}, \code{counterfactual-gate} \\
Decision science & \code{risk-economics-gate}, \code{uncertainty-calibration-gate},
\code{active-learning-gate}, \code{reviewer-sim-gate}, \code{explainable-ranking-gate} \\
Release \& interop & \code{rc-rehearsal-gate}, \code{leaderboard-gate}, \code{dossier-gate},
\code{interop-gate}, \code{artifact-mirror-gate} \\
Generalization & \code{framework-holdout-gate}, \code{perturbation-robustness-gate},
\code{historical-replay-study-gate}, \code{external-replication-kit-gate} \\
Trust stack & \code{signed-provenance-chain-gate}, \code{reproducible-build-attestation-gate},
\code{merkle-audit-log-gate}, \code{sbom-pinned-deps-gate}, \code{red-team-adversarial-gate} \\
Reviewer \& community & \code{quickstart-sixty-seconds-gate}, \code{inline-review-surface-gate},
\code{ci-pr-bot-gate}, \code{governance-policy-gate}, \code{citation-doi-gate} \\
Research frontier & \code{learned-risk-model-gate}, \code{neuro-symbolic-verdict-gate},
\code{theorem-prover-backend-gate}, \code{abstention-policy-gate}, \code{causal-effect-estimate-gate} \\
Camera-ready & \code{repro-appendix-gate}, \code{dataset-datasheet-gate}, \code{model-card-gate},
\code{evaluation-preregistration-gate}, \code{camera-ready-build-gate} \\
Mechanized foundations & \code{mechanized-operational-semantics-gate}, \code{soundness-theorem-gate},
\code{completeness-characterization-gate}, \code{refinement-types-columns-gate} \\
Ecosystem scale & \code{prevalence-estimator-gate}, \code{corpus-harness-gate},
\code{longitudinal-public-reruns-gate}, \code{drift-monitor-gate} \\
Adoption \& impact & \code{self-serve-onboarding-gate}, \code{policy-as-code-layer-gate},
\code{causal-effect-estimate-gate}, \code{grand-unified-evidence-index-gate} \\
\bottomrule
\caption{A representative selection from the 500+-gate catalog, grouped by theme.}
\label{tab:gatecatalog}
\end{longtable}

\section{Security and Privacy Model}
\label{app:security}

Because \tool{} ingests \emph{untrusted} repositories and archives, it carries an explicit threat
model that is itself gate-verified. Archive extraction defends against path traversal, symlink
escape, malformed archives, and archive bombs, while still extracting pinned public GitHub archives
through a content-addressed cache. Generated proposal files are non-executable, unapplied, and
excluded from native execution unless the user explicitly enables an allowlist of safe checks.
Redaction is byte-stable across repeated and resumed runs, secret-scanning proves that redacted
reports, prompts, bundles, and CI artifacts do not leak deterministic canary values, and an aggregate
privacy-metrics mode emits source-free risk trends with salted cohort identifiers so results can be
shared without raw evidence. Binaries, release archives, and public-corpus downloads carry
deterministic provenance and sorted checksums.

\section{Reproduction Guide}
\label{app:repro}

Every claim in this paper is regenerated by a \code{make} target. To reproduce the breadth results of
Section~\ref{sec:breadth} against pinned public repositories:

\begin{quote}\small\ttfamily
make multi-ecosystem-migration-gate \quad\# Laravel + 7-framework matrix\\
make nosql-change-gate \hspace{3.05em}\# Cassandra DROP on twitter-go + 5-engine matrix\\
make data-pipeline-gate \hspace{2.4em}\# Airflow+dbt+Spark on building-lakehouse\\
make schema-compat-gate \hspace{2.4em}\# Avro + protobuf on apache/avro\\
make infra-ordering-gate \hspace{2.0em}\# Kong/Anchore sequencing on helm/charts\\
make fixture-minimizer-gate \hspace{0.7em}\# 1-minimal reproductions per ecosystem\\
make parser-dashboard-gate \hspace{1.1em}\# coverage + fuzz robustness dashboard
\end{quote}

\noindent
Each target downloads only the pinned slice it needs, runs the real analyzer with \code{--no-llm},
asserts its invariants, and writes a \code{gate-summary.json}. A failing assertion---whether from a
code regression, a documentation drift, or a changed upstream ref---fails the target loudly, which is
precisely the property that makes the artifact trustworthy.

\end{document}
